\SourceLeaf{169}
\Needspace{9\baselineskip}
\pdfbookmark[1]{[IV.] Balance between the obligations and resources of the professors' widows' fund at Göttingen.}{witwen-iv}
\begin{center}
[IV.]\par\medskip
\textit{Balance between the obligations and resources of the professors' widows' fund at Göttingen.}
\end{center}
\noindent\hspace*{\parindent}In response to the wish conveyed to me some months ago by the university church deputation, I have undertaken a new calculation of the balance of the professors' widows' fund and have now completed this work. I shall arrange my report by first giving the necessary general explanations, then setting out the balance calculation itself in a concise, clearly arranged form, then discussing the individual items of the calculation more closely, and finally adding proposals for the period immediately ahead.
\begin{center}\rule{.26\textwidth}{.3pt}\end{center}
\noindent\hspace*{\parindent}The statement of the balance for an institution such as our widows' fund must evidently refer to a definite \textit{epoch}. My choice this time of 1 October 1851 will require no lengthy justification. The first balance calculation was made for 1 October 1845; for readily apparent reasons, the interval between two successive examinations must comprise a whole number of years to make the result stand out more clearly. Finally, if the new examination should make a change to the existing terms seem appropriate, the decision will evidently be based much more willingly on the latest state of affairs than on that of a year ago. The present statutes do prescribe a new revision every five years; but the time at which the request reached me no longer allowed another choice. Moreover, extending the five-year period to six years has caused no loss; rather, it has made the change in the state of affairs somewhat more distinct.

The essence of the entire balance calculation, both the present one and that of 1845, is that the institution's obligations on the one hand, and its resources on the other, are reduced to their equivalent capital value for the whole future, rather than for the next year or a few years. Only in this way is it possible to secure lasting viability for the whole future for an institution based only to a very small extent on contributions and chiefly on its assets, and whose membership has increased so greatly in recent decades.

Although this time, as in 1845, all calculations have been made twice, namely at an interest rate of \(3\tfrac{1}{2}\) and at one of 4 per cent, I have nevertheless considered it sufficient to give only the results at the former rate here. Part of the assets actually yields only \(3\tfrac{1}{2}\) per cent; for another part, now yielding a higher rate, a reduction in interest before too long is not improbable. In any case, prudence requires calculations of this kind always to be based on a rate somewhat below the one currently prevailing.

As mortality tables, in so far as the calculations depend on them, I have again used those of \textsc{Brune}, the most reliable in existence.

For the present balance calculation I have redetermined all those elements which can be derived only from the experience of the widows' fund itself, and which therefore gain in reliability when that experience covers a longer period, by combining the earlier experience with the new. Further details will be given under the individual items.

Finally, I note that all monetary amounts are to be understood in gold currency, and that although the original calculations were carried out throughout to fractions of a thaler, these fractions have been omitted in the present extract. This circumstance accounts for some apparent small
\SourceLeaf{170}
discrepancies in the stated sums, which may here and there amount to one or a couple of units.

\Needspace{210pt}
\begin{center}
\textit{Balance calculation of the widows\textquotesingle\ fund for 1 October 1851.}\par\smallskip
\footnotesize\setlength{\tabcolsep}{1pt}
\begin{minipage}[t]{.48\linewidth}
\centering
Liabilities.\par
\begin{tabular}{l r}
&Thalers\\
\shortstack[l]{Capital value of the fixed\\part of the pensions}&\\
\shortstack[l]{1. For the present\\widows . .}&28786\\
\shortstack[l]{2. For widows and orphans\\of the present members\\. . . . . . . . . . . . . . . . . . .}&58463\\
\shortstack[l]{3. For widows and orphans\\of members joining in future\\. . . . . . . . . . . .}&40712\\
\shortstack[l]{4. Capital value of the variable\\part of the pensions under\\the present rules . .}&25715\\
\shortstack[l]{5. Capital value of other\\expenditure . . .}&12907\\\cline{1-2}
\multicolumn{1}{r}{Total thalers}&166581
\end{tabular}
\end{minipage}\hfill\vrule\hfill
\begin{minipage}[t]{.48\linewidth}
\centering
Resources.\par
\begin{tabular}{l r}
&Thalers\\
Capital value&\\
\shortstack[l]{6. Of the pharmacy income\\. . . . . . . . .}&28571\\
\shortstack[l]{7. Of contributions from\\the present members . . .}&8961\\
\shortstack[l]{8. Of contributions from all\\members joining in future\\. . . . . . . . . . . . . . . . . . .}&13596\\
\shortstack[l]{9. Monetary assets of\\the widows\textquotesingle\ fund . . . . . .}&122290\\\cline{1-2}
\multicolumn{1}{r}{Total thalers}&173419
\end{tabular}\par\medskip
The balance calculation thus shows a \textit{surplus} of 6837 thalers.
\end{minipage}
\end{center}
